\normaltrue
\correctiontrue

%\UPSTIidClasse{11} % 11 sup, 12 spé
%\newcommand{\UPSTIidClasse}{12}

\exer{Mouvement TT -- $\star$ \label{CIN:02:B2:13:03}}
\setcounter{question}{0}\marginnote{\xpComp{CIN}{02}%\xpComp{CIN}{02}%\UPSTIcompetence[2]{C2-05}
\UPSTIcompetence[2]{B2-13}}
\index{Compétence C2-05}
\index{Compétence B2-13}\index{Compétence CIN-02}
\index{Mécanisme à 2 translations}
\ifcorrection
\else
\marginnote{\textbf{Pas de corrigé pour cet exercice.}}
\fi

\ifprof
\else
Soit le mécanisme suivant. On note $\vect{AB}=\lambda(t)\vect{i_0}$ et $\vect{BC}=\mu(t)\vect{j_0}$.
\begin{marginfigure}
\includegraphics[width=\linewidth]{03_TT_01}
\end{marginfigure}
\fi

% ==================
\ifprof
\else
\marginnote{
\begin{solution}
\begin{enumerate}
\item .
\item $x_C(t)= \lambda(t)$ et $y_C(t)= \mu(t)$.
\item $ \theta(t)=\dfrac{v}{R}t$.
\item $\lambda(t) = R\cos\left( \dfrac{v}{R}t\right)$, $\mu(t) = R\sin\left( \dfrac{v}{R}t\right)$.
\item .
\end{enumerate}
\end{solution}
Corrigé voir \ref{CIN:02:B2:13:03}}
\fi
% ==================

\question{Quel est le mouvement de \textbf{2} par rapport à~\textbf{0}.}
\ifprof ~\\
Le point $C$ a un mouvement quelconque dans le plan $\left(A,\vi{0},\vj{0}\right)$.
\else
\fi

\question{Donner l'équation du mouvement du point $C$ dans le mouvement de \textbf{2} par rapport à~\textbf{0}.}
\ifprof ~\\
On a $\vect{AC}=\lambda(t)\vect{i_0}+\mu(t)\vect{j_0}$ et donc
$\left\{
\begin{array}{l}
x_C(t)= \lambda(t) \\
y_C(t)= \mu(t)\\
z_C(t)= 0
\end{array}
\right.$
dans le repère $\repere{A}{i_0}{j_0}{k_0}$.
\else
\fi

On souhaite que le point $C$ réalise un cercle de centre $A$ et de rayon $R=\SI{10}{cm}$ à la vitesse $v=\SI{0,01}{m.s^{-1}}$.

\question{Donner la relation liant $\theta(t)$, $v$ et $R$.}

Par ailleurs la vitesse du point $C$ est donnée par $\vectv{C}{2}{0}=\deriv{\vect{AC}}{\rep{0}} = R\dot{\theta}\vect{e_\theta}$.
\ifprof ~\\
On a $v=R\dot{\theta}(t)$. Par intégration,
$\theta(t)=\dfrac{v}{R}t$ avec $\theta(0)=0$.
\else
\fi

\question{Donner les expressions de $\lambda(t)$ et $\mu(t)$ permettant la réalisation de cette trajectoire en fonction de $v$, $R$ et du temps.}
\ifprof ~\\
La trajectoire s'écrit
$\vect{AC}=R\vect{e_r}=R\cos\theta(t)\vi{0}+R\sin\theta(t)\vj{0}$.
Par identification,
$\lambda(t)=R\cos\theta(t)$ et $\mu(t)=R\sin\theta(t)$.

Au final,
$\left\{
\begin{array}{l}
\lambda(t)=R\cos\left(\dfrac{v}{R}t\right)\\
\mu(t)=R\sin\left(\dfrac{v}{R}t\right)
\end{array}
\right.$.
\else
\fi

\question{En utilisant Python, tracer $\lambda(t)$, $\mu(t)$ et la trajectoire générée.}
\ifprof
On choisit une durée correspondant à un tour complet :
\[
T=\frac{2\pi R}{v}.
\]
Sur l'intervalle $t\in[0,T]$, on évalue directement
\[
\lambda(t)=R\cos\left(\frac{v}{R}t\right),
\qquad
\mu(t)=R\sin\left(\frac{v}{R}t\right).
\]
Le premier tracé représente $\lambda$ et $\mu$ en fonction du temps. Le second représente la courbe paramétrée
\[
\bigl(x_C(t),y_C(t)\bigr)=\bigl(\lambda(t),\mu(t)\bigr),
\]
qui est un cercle de rayon $R$.

\begin{marginfigure}
\includegraphics[width=\linewidth]{03_TT_01_c}
\end{marginfigure}

\begin{marginfigure}
\includegraphics[width=\linewidth]{03_TT_02_c}
\end{marginfigure}
\else
\fi

% BEGIN AUTO CORRIGE INCLUDE
\InclureCorrige{corrige.tex}
% END AUTO CORRIGE INCLUDE
